\begin{table}[htp]
\centering

    % Esta primera línea define las columnas de la tabla. Los posibles tipos de columna son:
    % c: texto centrado
    % l: texto alineado a la izquierda
    % r: texto centrado a la derecha
    % p: columna de ancho fijo
    % Las columnas tienen ancho dinámico según la anchura máxima de los elementos que contengan.

    % Las columnas l/r/c no parten el texto en filas diferentes si éste es demasiado largo. Para ello, puede utilizar el tipo de columna de ancho fijo "p".
    
    % Las barras verticales | se usan para definir los bordes verticales de la tabla. Pruebe a eliminar algunas y observe qué ocurre.
    \begin{tabular}{ | p{3cm} | p{3cm} | p{3cm}| p{3cm} | }
        
        % A continuación van las filas de la tabla. En cada fila, las columnas se separan con el carácter &
        % Para terminar una fila se usa \\
        % Para incluir un borde horizontal entre filas se usa \hline

        % Cabecera con textos en negrita:
        \hline
        \textbf{Información} & \textbf{Funcionales} & \textbf{No funcionales} & \textbf{Hardware}\\
        \hline
        
        % Cuerpo de la tabla:
        I.01 Código GLSL & F.01 Biblioteca de shaders & NF.01 Robustez y detección de errores & H.01 Poseer una GPU configurada correctamente \\
        \hline
        I.02 Resultados, análisis y gráficas & F.02 Biblioteca de errores & NF.02 Modular y organizado & H.02 Sin mínima exigencia de CPU, RAM, disco, etc \\
        \hline
         & F.03 Biblioteca de polígonos & NF.03 Código comentado y en inglés & \\
        \hline
         & F.04 Biblioteca miscelánea de C & & \\
        \hline
         & F.05 Biblioteca de animación & & \\
        \hline
         & F.06 Ejemplo de uso y pruebas & & \\
        \hline
        
    \end{tabular} 
    
    \caption{Tabla de requisitos}
    \label{table:requirements} 
\end{table}
